% ============================================================
%  CHAPTER 2 — LITERATURE SURVEY
% ============================================================
\chapter{Literature Survey}

This chapter reviews existing research and literature in the fields of deep
learning-based facial expression recognition and human fatigue detection. A
summary of the most relevant studies is presented in the tables below,
followed by a brief analysis of how these works informed the development of
MindTrace.

% ── Survey Table Part 1 ──────────────────────────────────────
\begin{table}[H]
\centering
\caption{Literature Survey of Facial Emotion and Fatigue Detection (Part 1)}
\label{tab:lit_survey_1}
\renewcommand{\arraystretch}{1.5}
\begin{tabularx}{\textwidth}{|c|X|p{2.5cm}|X|p{2cm}|p{2cm}|}
\hline
\textbf{No.} & \textbf{Title \& Source} & \textbf{Techniques} & \textbf{Key Contribution} & \textbf{Pros} & \textbf{Cons} \\
\hline
1 & A Survey of Deep Learning FER Research -- IEEE Affective Computing, 2023 & CNN, ResNet, Attention & Comprehensive review of FER models (FER-2013, AffectNet). Discusses illumination and class imbalance. & Wide architecture coverage. & Lacks real-time integration. \\
\hline
2 & Improved FER Model Based on a Novel Deep Network -- Elsevier, 2024 & CNN, Batch Norm, Skip Conn. & Lightweight CNN with optimized residual links for high accuracy on RAF-DB. & Low latency; embedded ready. & Fails in poor lighting/ occlusion. \\
\hline
3 & Advancements in Driver Fatigue Detection -- Springer, 2024 & EAR, Facial Landmark Tracking & Hybrid system combining landmarks and blink pattern for early fatigue estimation. & Effective for safety systems. & Restricted to driver domain. \\
\hline
\end{tabularx}
\end{table}

\newpage

% ── Survey Table Part 2 ──────────────────────────────────────
\begin{table}[H]
\centering
\caption{Literature Survey of Facial Emotion and Fatigue Detection (Part 2)}
\label{tab:lit_survey_2}
\renewcommand{\arraystretch}{1.5}
\begin{tabularx}{\textwidth}{|c|X|p{2.5cm}|X|p{2cm}|p{2cm}|}
\hline
\textbf{No.} & \textbf{Title \& Source} & \textbf{Techniques} & \textbf{Key Contribution} & \textbf{Pros} & \textbf{Cons} \\
\hline
4 & Non-Invasive Fatigue Detection via Facial Parameters -- ACM HCI, 2025 & Emotion Tracking, Parameter Ext. & Fatigue detection using facial parameter variations for adaptive interaction. & User-friendly; non-contact. & Needs stable face position. \\
\hline
5 & Fatigue Monitoring Using Wearables and AI -- IEEE Sensors, 2025 & CNN + Sensor Fusion & Multimodal framework combining sensors (heart rate, blink) with CNN emotion recognition. & High precision and robustness. & Relies on external sensors. \\
\hline
\end{tabularx}
\end{table}

\section{Analysis of Existing Literature}

The reviewed literature highlights a significant transition in the field of
affective computing:

\begin{itemize}
    \item \textbf{From Pure Classification to Real-Time Systems}: Early
          works focused primarily on achieving high accuracy on static
          datasets. Recent trends, as seen in the work from Elsevier (2024),
          prioritize lightweight architectures suitable for deployment.
    
    \item \textbf{Introduction of Fatigue Metrics}: The integration of
          biometric signals like Eye Aspect Ratio (EAR), documented by
          Springer (2024), provides a deterministic way to measure focus
          and fatigue that complements emotional state analysis.
    
    \item \textbf{Multimodal vs.\ Visual-Only}: While multimodal systems
          (IEEE Sensors, 2025) offer high precision, the hardware
          complexity of wearables limits their adoption in standard
          e-learning environments. This justifies MindTrace's focus on
          purely visual, non-invasive monitoring.
\end{itemize}

MindTrace builds upon these foundations by combining the lightweight residual
learning approach for emotion classification with the non-invasive facial
parameter extraction for engagement scoring, delivering a unified real-time
system that operates on standard consumer hardware.
